\documentclass{article}
\usepackage[utf8]{inputenc}
\usepackage[russian]{babel} % For Russian language support
\usepackage{amsmath,amsfonts,amssymb} % For mathematical symbols
\usepackage{geometry}
\geometry{a4paper, margin=1in}

\title{Пример научной статьи}
\author{Автор}
\date{\today}

\begin{document}

\maketitle

\begin{abstract}
Данная статья представляет собой пример документа \LaTeX{} типа article с введением и заключением. В работе рассматриваются основные структурные элементы научной статьи.
\end{abstract}

\section{Введение}
Научные статьи играют важную роль в распространении знаний и результатов исследований. Формат \LaTeX{} позволяет создавать хорошо структурированные документы с высоким качеством типографики. Введение обычно содержит обзор темы, цели исследования и краткое изложение основных результатов.

В современном научном сообществе формат \LaTeX{} стал стандартом для подготовки статей, особенно в области математики, физики и информатики. Это связано с его возможностями по оформлению математических формул и библиографии.

\section{Основная часть}
Основная часть статьи содержит детальное изложение методов, результатов и анализа. В этой части обычно приводятся теоретические выкладки, экспериментальные данные и их интерпретация.

Пример математической формулы:
\begin{equation}
    E = mc^2
\end{equation}

Также в статье могут присутствовать таблицы, рисунки и ссылки на другие работы.

\section{Заключение}
В заключении подводятся итоги работы, формулируются основные выводы и обозначаются направления для дальнейших исследований. В этой статье мы рассмотрели основные структурные элементы научной статьи в формате \LaTeX{}.

Формат \LaTeX{} предоставляет мощные инструменты для создания профессионально выглядящих документов. Его использование позволяет сосредоточиться на содержании, а не на форматировании текста.

\end{document}